\section{Related Work}
\label{sec:related-work}

\subsection{Supply-Chain Attacks on Package Registries}
\label{sec:supply-chain}

Malicious publications in package registries are a well-documented,
persistent threat. Duan et al.~\cite{duan2021supply} measured hundreds of
malicious packages across PyPI, npm, and RubyGems, showing that most are
installed by ordinary users through ordinary \texttt{pip install}-style
commands, and that name-based imitation (typosquatting) is a dominant attack
style. Ohm et al.~\cite{ohm2020backstabber} cataloged attack techniques in
open-source software supply chains, including injections that trigger at
package-install time through \texttt{setup.py} hooks. Empirical studies of
the PyPI ecosystem report that malicious packages persist despite
registry-level safeguards~\cite{guo2023pypi, liang2021malicious}, and
graph-based and behavior-modeling approaches have been proposed for
detecting them in adjacent ecosystems~\cite{huang2024spiderscan}. Industry
efforts such as Datadog's malicious-software-packages
dataset~\cite{datadog-dataset} have made labeled collections of such
packages publicly available, which we use here.

\subsection{LLMs for Malicious-Code Detection}
\label{sec:llm-detection}

Recent work applies LLMs to security classification tasks. Studies on
vulnerability reasoning caution that LLMs remain unreliable on subtle
security judgments without auxiliary signal~\cite{ullah2024llmsec}. For
malicious packages specifically, LLMs have been used to detect malware in
the npm ecosystem~\cite{zahan2024npm}, and Ibiyo et al.~\cite{ease2025rag}---the closest
prior work to ours---evaluated LLMs with zero-shot prompting, RAG, and
few-shot learning on a dataset of malicious and benign PyPI
packages~\cite{ease2025rag}. Its knowledge bases were built from YARA rules,
GitHub security advisories, and malicious \texttt{setup.py} examples. A
central finding of that study was counterintuitive: retrieval augmentation
did not reliably help, yielding mediocre accuracy, whereas few-shot learning
was markedly more effective. The study, however, coupled retrieval to its
specific knowledge sources and a single retrieval mechanism, so it does not
separate the effect of \emph{whether} to retrieve from \emph{how} retrieval
is done.
A complementary line of work applies multi-agent LLM pipelines to
package-level evidence such as metadata and behavioral descriptions,
reporting improvements from aggregating multiple signals (\texttt{LAMPS},
``Many Hands Make Light Work'')~\cite{lamps2026}. Our work differs in scope
and design: instead of changing the knowledge source or inventing new
pipelines, we fix the corpus, the model, and the prompt, and vary only the
retrieval strategy, which lets us attribute performance differences to
retrieval design rather than to data or agent orchestration.

\subsection{Retrieval and Retrieval-Augmented Generation}
\label{sec:rag}

Retrieval-augmented generation combines a retriever with a generator so that
the model conditions on relevant external evidence~\cite{lewis2020rag,
gao2023ragsurvey}. Dense passage retrieval established embedding-based
retrieval as a standard for open-domain question answering~\cite{dpr2020};
in parallel, classical sparse scoring, notably the Okapi BM25 function,
remains a strong lexical baseline with a well-understood probabilistic
foundation~\cite{bm252009}. Hybrid systems fuse the two, most simply with
Reciprocal Rank Fusion~\cite{rrf2009}. Within security, RAG has been applied
to vulnerability detection in smart contracts~\cite{yu2024smartcontract} and
to knowledge-level augmentation of LLM vulnerability analysis with
CVE-derived documents~\cite{du2024vulrag}. For malicious code specifically,
retrieving code examples has been the first thing
tried~\cite{ease2025rag, crag2024}, yet---to our knowledge---no prior study
has systematically compared dense, sparse, hybrid, behavior-based, and
contrastive retrieval configurations on the same malicious-package benchmark
with the same generator. This is the gap we address.

\subsection{Summary}
\label{sec:related-summary}

In summary, prior work (i)~established the supply-chain threat and the
datasets we build on, (ii)~showed that LLMs augmented with retrieved
knowledge are promising but not sufficient for malicious-package detection,
and (iii)~provides the retrieval toolbox (dense, sparse, fusion). Our
contribution is a controlled, seven-condition comparison that isolates
retrieval strategy, together with faithfulness-oriented evaluation of the
model's behavior claims, on the same data and model family used by
Ibiyo et al.~\cite{ease2025rag}.
\endinput
